\solchapter{qinfo}

\sol{qinfo}{1} $\ket0$ is the north pole and $\ket1$ the south pole; $\ket+$ and $\ket-$ lie on the equator along $+x$ and $-x$, and $(\ket0+i\ket1)/\sqrt2$ along $+y$. Writing $H = (X+Z)/\sqrt2$, we get $H^2 = (X^2 + Z^2 + XZ + ZX)/2 = I$, since $X^2 = Z^2 = I$ and $XZ = -ZX$. Geometrically, $H$ is a half-turn about the axis halfway between $x$ and $z$: it swaps the $x$ and $z$ axes (so $\ket0\leftrightarrow\ket+$ and $\ket1\leftrightarrow\ket-$) and reverses $y$. Doing a half-turn twice gives the identity.

\sol{qinfo}{2} $\ket{\Phi^+} = (\ket{00}+\ket{11})/\sqrt2$ has no $\ket{01}$ or $\ket{10}$ component, so the results always agree. Substituting $\ket0 = (\ket++\ket-)/\sqrt2$ and $\ket1 = (\ket+-\ket-)/\sqrt2$, the cross terms cancel and $\ket{\Phi^+} = (\ket{++}+\ket{--})/\sqrt2$: the results agree in this basis too. No classical correlation does this in every basis at once.

\sol{qinfo}{3} $2^{50}\approx1.1\times10^{15}$ amplitudes, or $18$ petabytes --- the size of a large supercomputer's memory. For $300$ qubits, $2^{300}\approx2\times10^{90}$ amplitudes, more than the number of atoms in the observable universe ($\sim10^{80}$).

\sol{qinfo}{4} Take the inner product of $U\ket\phi\ket0 = \ket\phi\ket\phi$ with $U\ket\psi\ket0 = \ket\psi\ket\psi$. Since $U$ is unitary, the left side gives $\braket\psi\phi\braket00 = \braket\psi\phi$, and the right side gives $\braket\psi\phi^2$. So $x = x^2$ with $x = \braket\psi\phi$, which forces $x = 0$ or $1$: the machine can clone only states that are identical or orthogonal, i.e.\ a known set of classical alternatives.

\sol{qinfo}{5} $I\ket{\Phi^+} = \ket{\Phi^+}$; $X_1$ gives $(\ket{10}+\ket{01})/\sqrt2 = \ket{\Psi^+}$; $Z_1$ gives $(\ket{00}-\ket{11})/\sqrt2 = \ket{\Phi^-}$; and $X_1Z_1$ gives $(\ket{10}-\ket{01})/\sqrt2 = -\ket{\Psi^-}$. The four Bell states are orthogonal, so Bob can distinguish them with certainty: Alice's choice of one of four operations, two bits, arrives on a single qubit. (The other qubit, shared in advance, is the second ``bit''.)

\sol{qinfo}{6} Expand: $(\alpha\ket0+\beta\ket1)(\ket{00}+\ket{11})/\sqrt2 = (\alpha\ket{00}\ket0 + \alpha\ket{01}\ket1 + \beta\ket{10}\ket0 + \beta\ket{11}\ket1)/\sqrt2$, with Alice's two qubits written first. Substitute $\ket{00} = (\ket{\Phi^+}+\ket{\Phi^-})/\sqrt2$, $\ket{11} = (\ket{\Phi^+}-\ket{\Phi^-})/\sqrt2$, $\ket{01} = (\ket{\Psi^+}+\ket{\Psi^-})/\sqrt2$, and $\ket{10} = (\ket{\Psi^+}-\ket{\Psi^-})/\sqrt2$, and collect terms: this gives the stated expression. If Alice finds $\Phi^+$, Bob does nothing; $\Phi^-$, he applies $Z$; $\Psi^+$, he applies $X$; $\Psi^-$, he applies $X$ and then $Z$ (giving $-(\alpha\ket0+\beta\ket1)$, an irrelevant overall sign). Until Alice's two classical bits arrive, Bob's qubit is in an equal mixture of the four possibilities, which is maximally mixed and carries no information: nothing travels faster than light.

\sol{qinfo}{7} All four amplitudes start at $\tfrac12$. Flipping the marked one gives $(-\tfrac12,\tfrac12,\tfrac12,\tfrac12)$, with average $\tfrac14$. Reflecting about the average sends $a\mapsto2\cdot\tfrac14 - a$: the marked amplitude becomes $1$ and the others $0$. One query finds the item with certainty, where a classical search needs $2.25$ queries on average.

\sol{qinfo}{8} The correlation is $-\cos(\theta_a-\theta_b)$. With the given angles: $\expval{AB} = -\cos45^\circ$, $\expval{AB'} = -\cos45^\circ$, $\expval{A'B} = -\cos45^\circ$, and $\expval{A'B'} = -\cos135^\circ = +\cos45^\circ$. So $S = \abs{-4\cos45^\circ} = 2\sqrt2$. With $\vu b'$ at $135^\circ$, $\expval{AB'} = +\cos45^\circ$ and $\expval{A'B'} = -\cos45^\circ$, and the four terms cancel to $S = 0$. The signs in $S$ must line up with the geometry: each pair of axes used with a plus sign must be $45^\circ$ apart, and the pair with the minus sign $135^\circ$ apart.

\sol{qinfo}{9} $\rho_+ = \ket+\bra+ = \tfrac12\bigl(\begin{smallmatrix}1&1\\1&1\end{smallmatrix}\bigr)$ and $\rho_{\text{mix}} = \tfrac12\bigl(\begin{smallmatrix}1&0\\0&1\end{smallmatrix}\bigr)$. Both have diagonal entries $\tfrac12$, so a $Z$-measurement gives $0$ or $1$ with equal probability. For $X$, $\bra+\rho_+\ket+ = 1$ but $\bra+\rho_{\text{mix}}\ket+ = \tfrac12$: the pure state always gives $+$. The off-diagonal entries (the ``coherences'') are what distinguish a superposition from a mixture. $\rho_+^2 = \rho_+$ has trace $1$; $\rho_{\text{mix}}^2 = \tfrac14I$ has trace $\tfrac12$.

\sol{qinfo}{10} Alice and Bob keep only bits where they used the same basis. Eve guesses that basis right half the time; then she reads the bit correctly and resends it unchanged, causing no error. Half the time she guesses wrong; she then sends a state in the wrong basis, and Bob's measurement in the right one gives a random result, wrong with probability $\tfrac12$. The overall error rate is $\tfrac12\times\tfrac12 = 25\%$. Comparing a sample of their bits in public, Alice and Bob would spot this at once.

\sol{qinfo}{11} Both $\ket{000}$ and $\ket{111}$ have $Z_1Z_2 = Z_2Z_3 = +1$. A flip on qubit 1 makes both $\ket{100}$ and $\ket{011}$, which have $(Z_1Z_2, Z_2Z_3) = (-1,+1)$; similarly qubit 2 gives $(-1,-1)$ and qubit 3 gives $(+1,-1)$. The four syndromes are distinct, so they identify the error, which is then undone by applying the same $X$. Crucially, both branches of the superposition give the \emph{same} syndrome, so the state is an eigenstate of the measured observables and is not disturbed: the measurement reveals which qubit flipped but nothing about $\alpha$ or $\beta$. If two qubits flip, say 1 and 2, the syndrome $(+1,-1)$ blames qubit 3; ``correcting'' it gives $\alpha\ket{111}+\beta\ket{000}$, a logical error. A phase error $Z_1$ gives $\alpha\ket{000}-\beta\ket{111}$, which has the same syndrome as no error at all and is itself a logical error. Shor's nine-qubit code protects against both by encoding each of these three qubits again, in the $\ket\pm$ basis, where phase flips look like bit flips.

\sol{qinfo}{12} The box acting on $\ket x\ket- = \ket x(\ket0-\ket1)/\sqrt2$ gives $\ket x(\ket{f(x)}-\ket{1\oplus f(x)})/\sqrt2$, which is $\ket x\ket-$ if $f(x) = 0$ and $-\ket x\ket-$ if $f(x)=1$: so it is $(-1)^{f(x)}\ket x\ket-$. After the two Hadamards, the state is $\ket+\ket-$; after the box, the first qubit is
\[
  \tfrac{1}{\sqrt2}\big[(-1)^{f(0)}\ket0 + (-1)^{f(1)}\ket1\big],
\]
which is $\pm\ket+$ if $f(0) = f(1)$ and $\pm\ket-$ if not. The final Hadamard turns these into $\ket0$ and $\ket1$. One query answers a question about a global property of $f$ that classically needs two.

\sol{qinfo}{13} Write $C = A(B+B') + A'(B-B')$ and square, using $A^2 = A'^2 = B^2 = B'^2 = 1$ and the fact that Alice's operators commute with Bob's:
\[
  C^2 = (B+B')^2 + (B-B')^2 + AA'(B+B')(B-B') + A'A(B-B')(B+B').
\]
The first two terms are $(2 + BB'+B'B) + (2 - BB'-B'B) = 4$. Next, $(B+B')(B-B') = B'B - BB' = -[B,B']$ and $(B-B')(B+B') = [B,B']$, so the last two terms are $-(AA' - A'A)[B,B'] = -[A,A'][B,B']$. Now $\norm{[A,A']}\le\norm{AA'}+\norm{A'A}\le2$, and likewise for $B$, so $\norm{C^2}\le4+4 = 8$. Hence $\norm C\le2\sqrt2$ and $\abs{\expval C}\le2\sqrt2$. In a classical (local hidden variable) theory all the observables commute, the commutators vanish, $C^2 = 4$, and $\abs{\expval C}\le2$: Bell's inequality. The violation is made possible by non-commutativity, and limited by it.
